\documentclass{article}
\usepackage{amsmath}
\usepackage{graphicx}
\usepackage[colorlinks=true, allcolors=blue]{hyperref}

\title{Prestaciones y Ley de Amdhal}
\author{Victor Pinto}

\begin{document}
\maketitle



\section{Conceptos importantes}
El presente documento es un glosario de terminos y formulas que nos describen el concepto de prestaciones de un computador, metrica con la cual puede conocerse la rapidez con que la maquina realiza tareas, y poder comparar sus estadisticas con otros computadores.

1. Tiempo de Ejecucion: Se define como el tiempo total requerido por la computadora para completar una tarea, incluyendo las operaciones del disco duro, accesos a la memoria, la entrada y la salida, entre otros.

2. Ancho de banda o Productividad: Viene dado por el numero de tareas que un computador realiza por unidad de tiempo, esto es, en un periodo de tiempo definido.

3. Tiempo de CPU: Cantidad de tiempo que tarda el la CPU en terminar una tarea, sin incluir otras operaciones

4. Ciclos de Reloj: Tiempo de un periodo de reloj, el cual avanza a velocidad constante

5. Ciclos de Reloj por instruccion: Numero medio de ciclos de reloj que necesita una instruccion para ejecutarse (CPI)

6. Ley de Amdhal: Establece que el aumento de las prestaciones con una mejora determianda esta limitada por el uso que se le de a dicha mejora

\section{Formulas}
A continuacion se desglosan las formulas:

\[{Prestaciones}_{x} = \ \frac{1}{{Tiempo\ de\ Ejecucion}_{x}}\]

\[{Prestaciones}_{x} > \ {Prestaciones}_{y}\]

\[\frac{1}{{Tiempo\ de\ Ejecucion}_{x}} > \ \frac{1}{{Tiempo\ de\ Ejecucion}_{y}}\]

\[{Tiempo\ de\ Ejecucion}_{y} > \ {Tiempo\ de\ Ejecucion}_{x}\]

Para relacionar dos maquinas diferentes

\[\frac{{Prestaciones}_{x}}{{Prestaciones}_{y}} = n\]

Si X es n veces más rápida que Y, entonces el tiempo de ejecución de Y
es n veces

mayor que el de X

\(\frac{{Prestaciones}_{x}}{{Prestaciones}_{y}} = \ \frac{{Tiempo\ de\ Ejecucion}_{y}}{{Tiempo\ de\ Ejecucion}_{x}}\  = n\)

\[Tiempo\ \ de\ ejecucion\ de\ CPU\ para\ program.\  = \ Ciclos\ de\ reloj\ de\ la\ CPU*Tiemp.\ Ciclo\ Reloj\]

\[Ciclos\ de\ Reloj\ del\ CPU = Instrucciones\ de\ program*Media\ de\ ciclos\ por\ instrucc\]

\[Tiempo\ de\ Ejecucion = Num\ de\ instrucciones*CPI*Tiempo\ de\ Ciclo\]

\[Tiempo\ de\ Ejecucion = \frac{Num\ de\ instrucciones*CPI\ }{Frecuencia\ de\ reloj}\]

Ley de Amdhal :

\[Tiempo\ de\ ejecucion\ tras\ mejora = \ \frac{Tiempo\ de\ ejecucion\ por\ la\ mejora}{Cantidad\ de\ mejora} + Tiempo\ de\ ejec\ no\ afectado\]





\end{document}
